\documentclass[10pt,a4paper]{amsart}
\usepackage[margin=19mm]{geometry}
\usepackage{amsmath,amssymb,mathtools}
\usepackage[scr=rsfs,cal=euler]{mathalfa}
\usepackage{xfrac}
\usepackage[unicode,hidelinks,bookmarks=false]{hyperref}
\newcommand{\ie}{i.e.\ }
\newcommand{\cf}{cf.\ }
\newcommand{\resp}{respectively\ }
\newcommand{\Subsection}{\subsection}
\providecommand{\nobreakdash}{\nobreak\hbox{-}}
\setlength{\emergencystretch}{2em}
\multlinegap0pt
\usepackage{bm}
\usepackage{bbm}
\usepackage{dsfont}
\usepackage{xspace}
\usepackage{cancel}
\usepackage{mathtools}
\usepackage{amsthm}
\makeatletter
\@ifundefined{thm}{\newtheorem{thm}{Theorem}}{}
\@ifundefined{lemma}{\newtheorem{lemma}[thm]{Lemma}}{}
\makeatother
\newcommand{\abs}[1]{\left| #1 \right|}
\newcommand{\bracket}[1]{\left( #1 \right)}
\newcommand{\calx}{\mathcal{X}}
\title[Nonsimple closed geodesics with given intersection number on]{Nonsimple closed geodesics with given intersection number on hyperbolic surfaces}
\author{Wujie Shen}
\date{Source: arXiv:2305.00638v2 (CC BY 4.0)}
\begin{document}
\maketitle

\noindent\emph{Source and licence.} Nonsimple closed geodesics with given intersection number on hyperbolic surfaces. Version arXiv:2305.00638v2 (CC BY 4.0), distributed under the Creative Commons Attribution 4.0 International licence, as recorded in the arXiv licence metadata for this exact version and shown on the abstract page https://arxiv.org/abs/2305.00638v2. Full rights notice: licenses/arxiv-2305.00638-CC-BY-4.0.txt.\par\smallskip

\noindent\textit{Editorial scope.} Counted unit: the Thm whose source label is \texttt{m\_bigger2\_03} in the article (source0.tex lines 1608--1610, source label \texttt{m\_bigger2\_03}), delivered together with the article's own complete original proof of it (source0.tex lines 1612--1687, one \texttt{proof} environment of 76 lines). No mathematics is written, repaired or re-derived: statement and proof are the article's own text line for line. Required context retained uncounted: lemma\_32 (lines 1153--1165); m\_bigger2\_01 (lines 1451--1453). Every definition, display and label of those lines is retained unchanged. Omitted from this package and not redistributed: 1--1152, 1166--1450, 1454--1607, 1611--1611, 1688--2185. Nothing inside a retained excerpt is omitted or altered: every retained body is delivered exactly as the article's lines are written, so no omission receipt of any kind accompanies this package. Citations: the retained excerpts carry no citation command at all, so no bibliography entry of the article is transcribed and no reference is redistributed. Reference adaptation: none was needed and none was made; every reference inside the delivered unit resolves inside it, so no unresolved reference or citation is produced. Printed slips kept literal: no printed word, symbol, hypothesis or number of the counted statement or of its proof was altered. Layout and class: the article's own class is \texttt{article} with packages including amsmath, amsfonts, amscd, flafter, epsf, amssymb, epsfig, subfigure, bm, url, mathrsfs, amsthm, comment, geometry; the wrapper is the lane's pinned \texttt{amsart} editorial wrapper at 10pt on A4, which restates the theorem-like environments the retained text needs and adds the article's own macro definitions that the retained text needs (\texttt{\textbackslash abs}, \texttt{\textbackslash bracket}, \texttt{\textbackslash calx}), with extra wrapper packages (bm, bbm, dsfont, xspace, cancel, mathtools). The delivered numbering is the unit's own. Assets: no retained span contains a figure environment, a graphics-inclusion command, a PDF-page inclusion, a table or a TikZ picture; no image file is copied into this package, the package declares no asset dependency and no image file is redistributed with it. Licence: version arXiv:2305.00638v2 of this article is distributed under the Creative Commons Attribution 4.0 International licence, as recorded in the arXiv licence metadata for this exact version and shown on the abstract page https://arxiv.org/abs/2305.00638v2. A full rights notice is delivered at licenses/arxiv-2305.00638-CC-BY-4.0.txt. Characters: every retained excerpt is pure ASCII, so the delivered engine, XeLaTeX with its default Latin Modern text font, reports no missing character.
\begin{lemma}\label{lemma_32}

\begin{enumerate}
\item
   $w(\gamma_p)\leqslant 2\sinh\bracket{\frac{\ell(\gamma_p)}2}$. Similarly $w(\gamma_p')\leqslant \sinh\bracket{\frac{\ell(\gamma_p')}2}$.

\item
   For any $p$ such that $\gamma_p'\cap N_0(c_i)\neq\emptyset$ and $\gamma_p'$ is of general case, then $\ell(\gamma_p')\geqslant2\log(2+\sqrt3)$.

\item
For any $p$ such that $c_i\in\calx$, $\gamma_p\cap N_3(c_i)\neq\emptyset$, then $w(\gamma_p)\leqslant\frac{\ell(\gamma_p)}{\ell(c_i)}$.
\end{enumerate}
\end{lemma}

\begin{thm}\label{m_bigger2_01} 
 When $k>1750$ and $m\geqslant2$, if $m_2=0$ and $\abs{\Gamma\cap\Gamma}=k$, then $L\geqslant2\cosh^{-1}(2k+1)$. 
\end{thm}

\begin{thm}\label{m_bigger2_03} 
 When $k>1750$ and $m\geqslant2$, if $m_0,m_2\geqslant1$ and $\abs{\Gamma\cap\Gamma}=k$, then $L\geqslant2\cosh^{-1}(2k+1)$. 
\end{thm}

\begin{proof}
Recall that $\abs{\Gamma\cap\Gamma}=I(L_2',L_2'',m,m_1,m_2)$, where $L_1=L-L_2-L_2''$ and
\begin{align*}
I(L_2',L_2'',m,m_1,m_2)=\frac12\bracket{\frac{25}{12}L_1+m}^2+(m_1^2+2m_1m_2)e^{\frac{L_2'}{2m_1}}+{L_2''}e^{\frac{L_2''}4}+m^2-m
\end{align*}
For $a>0$ functions $f(x)=xe^{\frac{a}{x}}$ and $f(x)=x^2e^{\frac{a}{x}}$ are convex, hence $I(L_2',L_2'',m,m_1,m_2)$ is a convex function for $m_1\in(1,m_0)$. Hence one of the following holds:
\begin{align}\label{0001}
I\leqslant\frac12\bracket{\frac{25}{12}L_1+m}^2+(1+2m_2)e^{\frac{L_2'}{2}}+{L_2''}e^{\frac{L_2''}4}+m^2-m
\end{align}
\begin{align}\label{0002}
I\leqslant\frac12\bracket{\frac{25}{12}L_1+m}^2+(m_0^2+2m_0m_2)e^{\frac{L_2'}{2m_0}}+{L_2''}e^{\frac{L_2''}4}+m^2-m
\end{align}

\begin{enumerate}
\item
If (\ref{0001}) holds, since for $m_0+1\leqslant j\leqslant m$, $\ell(\gamma_j)\geqslant2\log2$, we have $L_2'\leqslant L_2-2m_2\log2\leqslant L-2m\log2-2m_2\log2$.  then when $m\geqslant7$ we have $m+m_2\geqslant8$, hence
\begin{align*}
 I&\leqslant\frac12\bracket{\frac{25}{12}L+m}^2+\frac{1+2m_2}{2^{m+m_2}}e^{\frac{L}{2}}+{L_2''}e^{\frac{L_2''}4}+m^2-m\\
&\leqslant\frac12\bracket{\frac{25}{12}L+m}^2+\frac{3}{256}e^{\frac{L}{2}}+(L-2m\log2)e^{\frac{L-2m\log2}4}+m^2-m\\
&\leqslant\max\left\{\frac12\bracket{\frac{25}{12}L+7}^2+\frac{3}{256}e^{\frac{L}{2}}+Le^{\frac{L-14\log2}4}+42,\frac12\bracket{\frac{25}{12}L+\frac{L}{2\log2}}^2+\frac{3}{256}e^{\frac{L}{2}}+\bracket{\frac{L}{2\log2}}^2\right\}\\
&\leqslant\max\left\{\frac12\bracket{\frac{25}{12}L+7}^2+\frac{3}{256}e^{\frac{L}{2}}+\frac{L}{8\sqrt2}e^{\frac{L}4}+42,\frac{3}{256}e^{\frac{L}{2}}+4.46L^2\right\}
\end{align*}
The second inequality uses the fact that when $m>m_2\geqslant1$ and $m\geqslant7$, then $\frac{1+2m_2}{2^{m+m_2}}\leqslant\frac3{256}$. The third inequality uses the convexity on $m\in(7,\frac{L}{2\log2})$. When $L>17.7$ we have $4I+2\leqslant e^{\frac{L}2}$, and when $L\leqslant17.7$ we have $4.46L^2+\frac{3}{256}e^{\frac{L}{2}}<1700$ and $\frac12\bracket{\frac{25}{12}L+7}^2+\frac{3}{256}e^{\frac{L}{2}}+\frac{L}{8\sqrt2}e^{\frac{L}4}+42<1700$, contradiction.

When $4\leqslant m\leqslant6$, using the convexity of $L_2\in(0,L-2m\log2)$, and $L_2\leqslant L-2m\log2\leqslant L-8\log2$, we have
\begin{align*}
 I&\leqslant\frac12\bracket{\frac{25}{12}(L-L_2)+m}^2+(1+2m_2)e^{\frac{L_2-2m_2\log2}2}+{L_2}e^{\frac{L_2}4}+m^2-m\\
&\leqslant\max\left\{m^2-m+2m_2+1+\frac12\bracket{\frac{25}{12}L+6}^2, \frac12\bracket{\frac{25\log2}{6}m+m}^2+\frac{3}{32}e^{\frac{L}{2}}+\frac{L}{4}e^{\frac{L}4}+m^2-m\right\}\\
&\leqslant\max\left\{44+\frac12\bracket{\frac{25}{12}L+6}^2, 303+\frac{3}{32}e^{\frac{L}{2}}+\frac{L}{4}e^{\frac{L}4}\right\}
\end{align*}
since $\frac{1+2m_2}{2^{m+m_2}}\leqslant\frac3{32}$ when $m\geqslant4, m_2\geqslant1$. Hence when $L>17.7$ then $4I+2\leqslant e^{\frac{L}2}$, when $L\leqslant17.7$ then $I\leqslant1750$, contradiction.

When $m=3$, $m_1\geqslant1$ implies $\ell(\gamma_1')\geqslant2\log(2+\sqrt3)>2.62$ using Lemma \ref{lemma_32}, hence $L_2''\leqslant L-4\log2-2\log(2+\sqrt3)\in(L-5.41,L-5.4)$. Using the convexity of $L_2\in(0,L-6\log2)$, we have
\begin{align*}
 I&\leqslant\frac12\bracket{\frac{25}{12}(L-L_2)+3}^2+(1+2m_2)e^{\frac{L_2-2m_2\log2}{2}}+{L_2}e^{\frac{L-4\log2-2\log(2+\sqrt3)}4}+6\\
&\leqslant\max\left\{13+\frac12\bracket{\frac{25}{12}L+6}^2,\frac12\bracket{\frac{25\times6\log2}{12}+3}^2+\frac{3}{16}e^{\frac{L}{2}}+\frac{L-6\log2}{2\sqrt{2+\sqrt3}}e^{\frac{L}4}+6\right\}\\
&\leqslant\max\left\{13+\frac12\bracket{\frac{25}{12}L+6}^2, 110+\frac{3}{16}e^{\frac{L}{2}}+\frac{L-6\log2}{3.86}e^{\frac{L}4}\right\}
\end{align*}
since $\frac{1+2m_2}{2^{m+m_2}}\leqslant\frac3{16}$ when $m\geqslant3, m_2\geqslant1$. Hence when $L>17.7$ then $4I+2\leqslant e^{\frac{L}2}$, when $L\leqslant17.7$ then $I\leqslant1750$, contradiction.




When $m_1=m_2=1$, then suppose $\gamma_1$ is of general case and $\gamma_2$ is of special case, and $\gamma_1\cap N_0(c_i)\neq\emptyset$, $\gamma_2\cap N_0(c_{i'})\neq\emptyset$. If $i\neq i'$ then $\abs{\gamma_2\cap\gamma_2}=\abs{\gamma_1\cap\gamma_2}=0$, same as Theorem \ref{m_bigger2_01} to get the proof. Otherwise $i=i'$. 

If $\ell(c_i)\geqslant 0.25$ then $w(\gamma_1)\leqslant\frac{L_2}{0.25}<4L$, and $\abs{\Gamma_2\cap\Gamma_2}\leqslant2w(\gamma_1)+4$, hence 
$$k=\abs{\Gamma\cap\Gamma}\leqslant\frac12\bracket{\frac{25}{12}L+2}^2+8L+4$$ 
When $L>17.7$ we have $4(\frac12\bracket{\frac{25}{12}L+2}^2+8L+4)+2<e^{\frac{L}2}$, when $L\leqslant17.7$ we have $k<1750$, ontradiction.

If $\ell(c_i)<0.25$, we have $\ell(\gamma_2)\geqslant6\log2$, hence $\ell(\gamma_1)\leqslant L-10\log2$. Hence using Lemma \ref{lemma_32} we have $w(\gamma_1)\leqslant\frac1{32}e^{\frac{L}2}$, then
$$k=\abs{\Gamma\cap\Gamma}\leqslant\frac12\bracket{\frac{25}{12}L+2}^2+\frac1{16}e^{\frac{L}2}+4$$
When $L>17.7$ we have $4k+2<e^{\frac{L}2}$, when $L\leqslant17.7$ we have $k<1750$, ontradiction.


\item
If (\ref{0002}) holds and $m_0\geqslant2$, then $\frac{L_2'}{2m_0}\leqslant\frac{L_2'}4$, hence when $L\geqslant16.7$ we have
\begin{align*}
I&\leqslant\frac12\bracket{\frac{25}{12}L_1+m}^2+m^2e^{\frac{L_2'}{2m_0}}+{L_2''}e^{\frac{L_2''}4}+m^2-m\\
&\leqslant\frac12\bracket{\frac{25}{12}L+m}^2+\frac{m^2+L}{2^{\frac{m}2}}e^{\frac{L}4}+m^2-m\\
&\leqslant\max\left\{\frac12\bracket{\frac{25}{12}L+2}^2+(2+\frac{L}2)e^{\frac{L}4}+2, \frac12\bracket{\frac{25}{12}L+\frac{L}{2\log2}}^2+2\bracket{\frac{L}{2\log2}}^2+L\right\}
\end{align*}
The third inequality using the convexity on $m\in(2,\frac{L}{2\log2})$, here when $L\geqslant16.7$ the function $f(x)=\frac{x^2+L}{2^{\frac{x}2}}$ is convex for $x\in(2,+\infty)$. When $L\geqslant17.7$ then $4I+2\leqslant e^{\frac{L}2}$, and when $16.7\leqslant L<17.7$, $I<1750$, contradiction.

When $L<16.7$, since for $x>1$, $\frac{x^2}{2^{\frac{x}2}}<4.51$, then using the convexity on $m\in(2,\frac{L}{2\log2})$ we have
\begin{align*}
I&\leqslant\frac12\bracket{\frac{25}{12}L_1+m}^2+(4.51+\frac{L}{2^{\frac{m}2}})e^{\frac{L}4}+m^2-m\\
&\leqslant\max\left\{\frac12\bracket{\frac{25}{12}L+2}^2+(4.51+\frac{L}2)e^{\frac{L}4}+2, \frac12\bracket{\frac{25}{12}L+\frac{L}{2\log2}}^2+\bracket{\frac{L}{2\log2}}^2+L+4.51e^{\frac{L}4}\right\}\\
&<1750
\end{align*}
a contradiction.

\end{enumerate}



\end{proof}

\end{document}
